\section{}

\paragraph{1145年}\eventanchor{XIA-1145-REGNAL-YEAR}{XIA-1145-REGNAL-YEAR}　西夏以李仁孝在位第6年作为诏令、赋役与外交文书的纪年框架。账本以《宋史》卷485〈夏国传〉；《辽史》《金史》相关本纪；黑水城西夏文文书与考古报告标明来源路径；该锚点连接编年节点，关于数量、实施和地方后果仍应遵守材料的文体与立场限制。

\paragraph{1146年}\eventanchor{XIA-1146-REGNAL-YEAR}{XIA-1146-REGNAL-YEAR}　西夏以李仁孝在位第7年作为诏令、赋役与外交文书的纪年框架。账本以《宋史》卷485〈夏国传〉；《辽史》《金史》相关本纪；黑水城西夏文文书与考古报告标明来源路径；该锚点连接编年节点，关于数量、实施和地方后果仍应遵守材料的文体与立场限制。

\paragraph{1147年}\eventanchor{XIA-1147-REGNAL-YEAR}{XIA-1147-REGNAL-YEAR}　西夏以李仁孝在位第8年作为诏令、赋役与外交文书的纪年框架。账本以《宋史》卷485〈夏国传〉；《辽史》《金史》相关本纪；黑水城西夏文文书与考古报告标明来源路径；该锚点连接编年节点，关于数量、实施和地方后果仍应遵守材料的文体与立场限制。

\paragraph{1148年}\eventanchor{XIA-1148-REGNAL-YEAR}{XIA-1148-REGNAL-YEAR}　西夏以李仁孝在位第9年作为诏令、赋役与外交文书的纪年框架。账本以《宋史》卷485〈夏国传〉；《辽史》《金史》相关本纪；黑水城西夏文文书与考古报告标明来源路径；该锚点连接编年节点，关于数量、实施和地方后果仍应遵守材料的文体与立场限制。

\paragraph{1149年}\eventanchor{XIA-1149-REGNAL-YEAR}{XIA-1149-REGNAL-YEAR}　西夏以李仁孝在位第10年作为诏令、赋役与外交文书的纪年框架。账本以《宋史》卷485〈夏国传〉；《辽史》《金史》相关本纪；黑水城西夏文文书与考古报告标明来源路径；该锚点连接编年节点，关于数量、实施和地方后果仍应遵守材料的文体与立场限制。

\paragraph{1150年}\eventanchor{XIA-1150-REGNAL-YEAR}{XIA-1150-REGNAL-YEAR}　西夏以李仁孝在位第11年作为诏令、赋役与外交文书的纪年框架。账本以《宋史》卷485〈夏国传〉；《辽史》《金史》相关本纪；黑水城西夏文文书与考古报告标明来源路径；该锚点连接编年节点，关于数量、实施和地方后果仍应遵守材料的文体与立场限制。

\paragraph{1151年}\eventanchor{XIA-1151-REGNAL-YEAR}{XIA-1151-REGNAL-YEAR}　西夏以李仁孝在位第12年作为诏令、赋役与外交文书的纪年框架。账本以《宋史》卷485〈夏国传〉；《辽史》《金史》相关本纪；黑水城西夏文文书与考古报告标明来源路径；该锚点连接编年节点，关于数量、实施和地方后果仍应遵守材料的文体与立场限制。

\paragraph{1152年}\eventanchor{XIA-1152-REGNAL-YEAR}{XIA-1152-REGNAL-YEAR}　西夏以李仁孝在位第13年作为诏令、赋役与外交文书的纪年框架。账本以《宋史》卷485〈夏国传〉；《辽史》《金史》相关本纪；黑水城西夏文文书与考古报告标明来源路径；该锚点连接编年节点，关于数量、实施和地方后果仍应遵守材料的文体与立场限制。

\paragraph{1153年}\eventanchor{XIA-1153-REGNAL-YEAR}{XIA-1153-REGNAL-YEAR}　西夏以李仁孝在位第14年作为诏令、赋役与外交文书的纪年框架。账本以《宋史》卷485〈夏国传〉；《辽史》《金史》相关本纪；黑水城西夏文文书与考古报告标明来源路径；该锚点连接编年节点，关于数量、实施和地方后果仍应遵守材料的文体与立场限制。

\paragraph{1154年}\eventanchor{XIA-1154-REGNAL-YEAR}{XIA-1154-REGNAL-YEAR}　西夏以李仁孝在位第15年作为诏令、赋役与外交文书的纪年框架。账本以《宋史》卷485〈夏国传〉；《辽史》《金史》相关本纪；黑水城西夏文文书与考古报告标明来源路径；该锚点连接编年节点，关于数量、实施和地方后果仍应遵守材料的文体与立场限制。

\paragraph{1155年}\eventanchor{XIA-1155-REGNAL-YEAR}{XIA-1155-REGNAL-YEAR}　西夏以李仁孝在位第16年作为诏令、赋役与外交文书的纪年框架。账本以《宋史》卷485〈夏国传〉；《辽史》《金史》相关本纪；黑水城西夏文文书与考古报告标明来源路径；该锚点连接编年节点，关于数量、实施和地方后果仍应遵守材料的文体与立场限制。

\paragraph{1156年}\eventanchor{XIA-1156-REGNAL-YEAR}{XIA-1156-REGNAL-YEAR}　西夏以李仁孝在位第17年作为诏令、赋役与外交文书的纪年框架。账本以《宋史》卷485〈夏国传〉；《辽史》《金史》相关本纪；黑水城西夏文文书与考古报告标明来源路径；该锚点连接编年节点，关于数量、实施和地方后果仍应遵守材料的文体与立场限制。

\paragraph{1157年}\eventanchor{XIA-1157-REGNAL-YEAR}{XIA-1157-REGNAL-YEAR}　西夏以李仁孝在位第18年作为诏令、赋役与外交文书的纪年框架。账本以《宋史》卷485〈夏国传〉；《辽史》《金史》相关本纪；黑水城西夏文文书与考古报告标明来源路径；该锚点连接编年节点，关于数量、实施和地方后果仍应遵守材料的文体与立场限制。

\paragraph{1158年}\eventanchor{XIA-1158-REGNAL-YEAR}{XIA-1158-REGNAL-YEAR}　西夏以李仁孝在位第19年作为诏令、赋役与外交文书的纪年框架。账本以《宋史》卷485〈夏国传〉；《辽史》《金史》相关本纪；黑水城西夏文文书与考古报告标明来源路径；该锚点连接编年节点，关于数量、实施和地方后果仍应遵守材料的文体与立场限制。

\paragraph{1159年}\eventanchor{XIA-1159-REGNAL-YEAR}{XIA-1159-REGNAL-YEAR}　西夏以李仁孝在位第20年作为诏令、赋役与外交文书的纪年框架。账本以《宋史》卷485〈夏国传〉；《辽史》《金史》相关本纪；黑水城西夏文文书与考古报告标明来源路径；该锚点连接编年节点，关于数量、实施和地方后果仍应遵守材料的文体与立场限制。

\paragraph{1160年}\eventanchor{XIA-1160-REGNAL-YEAR}{XIA-1160-REGNAL-YEAR}　西夏以李仁孝在位第21年作为诏令、赋役与外交文书的纪年框架。账本以《宋史》卷485〈夏国传〉；《辽史》《金史》相关本纪；黑水城西夏文文书与考古报告标明来源路径；该锚点连接编年节点，关于数量、实施和地方后果仍应遵守材料的文体与立场限制。

\paragraph{1161年}\eventanchor{XIA-1161-REGNAL-YEAR}{XIA-1161-REGNAL-YEAR}　西夏以李仁孝在位第22年作为诏令、赋役与外交文书的纪年框架。账本以《宋史》卷485〈夏国传〉；《辽史》《金史》相关本纪；黑水城西夏文文书与考古报告标明来源路径；该锚点连接编年节点，关于数量、实施和地方后果仍应遵守材料的文体与立场限制。

\paragraph{1162年}\eventanchor{XIA-1162-REGNAL-YEAR}{XIA-1162-REGNAL-YEAR}　西夏以李仁孝在位第23年作为诏令、赋役与外交文书的纪年框架。账本以《宋史》卷485〈夏国传〉；《辽史》《金史》相关本纪；黑水城西夏文文书与考古报告标明来源路径；该锚点连接编年节点，关于数量、实施和地方后果仍应遵守材料的文体与立场限制。

\paragraph{1163年}\eventanchor{XIA-1163-REGNAL-YEAR}{XIA-1163-REGNAL-YEAR}　西夏以李仁孝在位第24年作为诏令、赋役与外交文书的纪年框架。账本以《宋史》卷485〈夏国传〉；《辽史》《金史》相关本纪；黑水城西夏文文书与考古报告标明来源路径；该锚点连接编年节点，关于数量、实施和地方后果仍应遵守材料的文体与立场限制。

\paragraph{1164年}\eventanchor{XIA-1164-REGNAL-YEAR}{XIA-1164-REGNAL-YEAR}　西夏以李仁孝在位第25年作为诏令、赋役与外交文书的纪年框架。账本以《宋史》卷485〈夏国传〉；《辽史》《金史》相关本纪；黑水城西夏文文书与考古报告标明来源路径；该锚点连接编年节点，关于数量、实施和地方后果仍应遵守材料的文体与立场限制。

\paragraph{1165年}\eventanchor{XIA-1165-REGNAL-YEAR}{XIA-1165-REGNAL-YEAR}　西夏以李仁孝在位第26年作为诏令、赋役与外交文书的纪年框架。账本以《宋史》卷485〈夏国传〉；《辽史》《金史》相关本纪；黑水城西夏文文书与考古报告标明来源路径；该锚点连接编年节点，关于数量、实施和地方后果仍应遵守材料的文体与立场限制。

\paragraph{1166年}\eventanchor{XIA-1166-REGNAL-YEAR}{XIA-1166-REGNAL-YEAR}　西夏以李仁孝在位第27年作为诏令、赋役与外交文书的纪年框架。账本以《宋史》卷485〈夏国传〉；《辽史》《金史》相关本纪；黑水城西夏文文书与考古报告标明来源路径；该锚点连接编年节点，关于数量、实施和地方后果仍应遵守材料的文体与立场限制。

\paragraph{1167年}\eventanchor{XIA-1167-REGNAL-YEAR}{XIA-1167-REGNAL-YEAR}　西夏以李仁孝在位第28年作为诏令、赋役与外交文书的纪年框架。账本以《宋史》卷485〈夏国传〉；《辽史》《金史》相关本纪；黑水城西夏文文书与考古报告标明来源路径；该锚点连接编年节点，关于数量、实施和地方后果仍应遵守材料的文体与立场限制。

\paragraph{1168年}\eventanchor{XIA-1168-REGNAL-YEAR}{XIA-1168-REGNAL-YEAR}　西夏以李仁孝在位第29年作为诏令、赋役与外交文书的纪年框架。账本以《宋史》卷485〈夏国传〉；《辽史》《金史》相关本纪；黑水城西夏文文书与考古报告标明来源路径；该锚点连接编年节点，关于数量、实施和地方后果仍应遵守材料的文体与立场限制。

\paragraph{1169年}\eventanchor{XIA-1169-REGNAL-YEAR}{XIA-1169-REGNAL-YEAR}　西夏以李仁孝在位第30年作为诏令、赋役与外交文书的纪年框架。账本以《宋史》卷485〈夏国传〉；《辽史》《金史》相关本纪；黑水城西夏文文书与考古报告标明来源路径；该锚点连接编年节点，关于数量、实施和地方后果仍应遵守材料的文体与立场限制。

\paragraph{1170年}\eventanchor{XIA-1170-REGNAL-YEAR}{XIA-1170-REGNAL-YEAR}　西夏以李仁孝在位第31年作为诏令、赋役与外交文书的纪年框架。账本以《宋史》卷485〈夏国传〉；《辽史》《金史》相关本纪；黑水城西夏文文书与考古报告标明来源路径；该锚点连接编年节点，关于数量、实施和地方后果仍应遵守材料的文体与立场限制。

\paragraph{1171年}\eventanchor{XIA-1171-REGNAL-YEAR}{XIA-1171-REGNAL-YEAR}　西夏以李仁孝在位第32年作为诏令、赋役与外交文书的纪年框架。账本以《宋史》卷485〈夏国传〉；《辽史》《金史》相关本纪；黑水城西夏文文书与考古报告标明来源路径；该锚点连接编年节点，关于数量、实施和地方后果仍应遵守材料的文体与立场限制。

\paragraph{1172年}\eventanchor{XIA-1172-REGNAL-YEAR}{XIA-1172-REGNAL-YEAR}　西夏以李仁孝在位第33年作为诏令、赋役与外交文书的纪年框架。账本以《宋史》卷485〈夏国传〉；《辽史》《金史》相关本纪；黑水城西夏文文书与考古报告标明来源路径；该锚点连接编年节点，关于数量、实施和地方后果仍应遵守材料的文体与立场限制。

\paragraph{1173年}\eventanchor{XIA-1173-REGNAL-YEAR}{XIA-1173-REGNAL-YEAR}　西夏以李仁孝在位第34年作为诏令、赋役与外交文书的纪年框架。账本以《宋史》卷485〈夏国传〉；《辽史》《金史》相关本纪；黑水城西夏文文书与考古报告标明来源路径；该锚点连接编年节点，关于数量、实施和地方后果仍应遵守材料的文体与立场限制。

\paragraph{1174年}\eventanchor{XIA-1174-REGNAL-YEAR}{XIA-1174-REGNAL-YEAR}　西夏以李仁孝在位第35年作为诏令、赋役与外交文书的纪年框架。账本以《宋史》卷485〈夏国传〉；《辽史》《金史》相关本纪；黑水城西夏文文书与考古报告标明来源路径；该锚点连接编年节点，关于数量、实施和地方后果仍应遵守材料的文体与立场限制。

\paragraph{1175年}\eventanchor{XIA-1175-REGNAL-YEAR}{XIA-1175-REGNAL-YEAR}　西夏以李仁孝在位第36年作为诏令、赋役与外交文书的纪年框架。账本以《宋史》卷485〈夏国传〉；《辽史》《金史》相关本纪；黑水城西夏文文书与考古报告标明来源路径；该锚点连接编年节点，关于数量、实施和地方后果仍应遵守材料的文体与立场限制。

\paragraph{1176年}\eventanchor{XIA-1176-REGNAL-YEAR}{XIA-1176-REGNAL-YEAR}　西夏以李仁孝在位第37年作为诏令、赋役与外交文书的纪年框架。账本以《宋史》卷485〈夏国传〉；《辽史》《金史》相关本纪；黑水城西夏文文书与考古报告标明来源路径；该锚点连接编年节点，关于数量、实施和地方后果仍应遵守材料的文体与立场限制。

\paragraph{1177年}\eventanchor{XIA-1177-REGNAL-YEAR}{XIA-1177-REGNAL-YEAR}　西夏以李仁孝在位第38年作为诏令、赋役与外交文书的纪年框架。账本以《宋史》卷485〈夏国传〉；《辽史》《金史》相关本纪；黑水城西夏文文书与考古报告标明来源路径；该锚点连接编年节点，关于数量、实施和地方后果仍应遵守材料的文体与立场限制。

\paragraph{1178年}\eventanchor{XIA-1178-REGNAL-YEAR}{XIA-1178-REGNAL-YEAR}　西夏以李仁孝在位第39年作为诏令、赋役与外交文书的纪年框架。账本以《宋史》卷485〈夏国传〉；《辽史》《金史》相关本纪；黑水城西夏文文书与考古报告标明来源路径；该锚点连接编年节点，关于数量、实施和地方后果仍应遵守材料的文体与立场限制。

\paragraph{1179年}\eventanchor{XIA-1179-REGNAL-YEAR}{XIA-1179-REGNAL-YEAR}　西夏以李仁孝在位第40年作为诏令、赋役与外交文书的纪年框架。账本以《宋史》卷485〈夏国传〉；《辽史》《金史》相关本纪；黑水城西夏文文书与考古报告标明来源路径；该锚点连接编年节点，关于数量、实施和地方后果仍应遵守材料的文体与立场限制。

\paragraph{1180年}\eventanchor{XIA-1180-REGNAL-YEAR}{XIA-1180-REGNAL-YEAR}　西夏以李仁孝在位第41年作为诏令、赋役与外交文书的纪年框架。账本以《宋史》卷485〈夏国传〉；《辽史》《金史》相关本纪；黑水城西夏文文书与考古报告标明来源路径；该锚点连接编年节点，关于数量、实施和地方后果仍应遵守材料的文体与立场限制。

\paragraph{1181年}\eventanchor{XIA-1181-REGNAL-YEAR}{XIA-1181-REGNAL-YEAR}　西夏以李仁孝在位第42年作为诏令、赋役与外交文书的纪年框架。账本以《宋史》卷485〈夏国传〉；《辽史》《金史》相关本纪；黑水城西夏文文书与考古报告标明来源路径；该锚点连接编年节点，关于数量、实施和地方后果仍应遵守材料的文体与立场限制。

\paragraph{1182年}\eventanchor{XIA-1182-REGNAL-YEAR}{XIA-1182-REGNAL-YEAR}　西夏以李仁孝在位第43年作为诏令、赋役与外交文书的纪年框架。账本以《宋史》卷485〈夏国传〉；《辽史》《金史》相关本纪；黑水城西夏文文书与考古报告标明来源路径；该锚点连接编年节点，关于数量、实施和地方后果仍应遵守材料的文体与立场限制。

\paragraph{1183年}\eventanchor{XIA-1183-REGNAL-YEAR}{XIA-1183-REGNAL-YEAR}　西夏以李仁孝在位第44年作为诏令、赋役与外交文书的纪年框架。账本以《宋史》卷485〈夏国传〉；《辽史》《金史》相关本纪；黑水城西夏文文书与考古报告标明来源路径；该锚点连接编年节点，关于数量、实施和地方后果仍应遵守材料的文体与立场限制。

\paragraph{1184年}\eventanchor{XIA-1184-REGNAL-YEAR}{XIA-1184-REGNAL-YEAR}　西夏以李仁孝在位第45年作为诏令、赋役与外交文书的纪年框架。账本以《宋史》卷485〈夏国传〉；《辽史》《金史》相关本纪；黑水城西夏文文书与考古报告标明来源路径；该锚点连接编年节点，关于数量、实施和地方后果仍应遵守材料的文体与立场限制。

\paragraph{1185年}\eventanchor{XIA-1185-REGNAL-YEAR}{XIA-1185-REGNAL-YEAR}　西夏以李仁孝在位第46年作为诏令、赋役与外交文书的纪年框架。账本以《宋史》卷485〈夏国传〉；《辽史》《金史》相关本纪；黑水城西夏文文书与考古报告标明来源路径；该锚点连接编年节点，关于数量、实施和地方后果仍应遵守材料的文体与立场限制。

\paragraph{1186年}\eventanchor{XIA-1186-REGNAL-YEAR}{XIA-1186-REGNAL-YEAR}　西夏以李仁孝在位第47年作为诏令、赋役与外交文书的纪年框架。账本以《宋史》卷485〈夏国传〉；《辽史》《金史》相关本纪；黑水城西夏文文书与考古报告标明来源路径；该锚点连接编年节点，关于数量、实施和地方后果仍应遵守材料的文体与立场限制。

\paragraph{1187年}\eventanchor{XIA-1187-REGNAL-YEAR}{XIA-1187-REGNAL-YEAR}　西夏以李仁孝在位第48年作为诏令、赋役与外交文书的纪年框架。账本以《宋史》卷485〈夏国传〉；《辽史》《金史》相关本纪；黑水城西夏文文书与考古报告标明来源路径；该锚点连接编年节点，关于数量、实施和地方后果仍应遵守材料的文体与立场限制。

\paragraph{1188年}\eventanchor{XIA-1188-REGNAL-YEAR}{XIA-1188-REGNAL-YEAR}　西夏以李仁孝在位第49年作为诏令、赋役与外交文书的纪年框架。账本以《宋史》卷485〈夏国传〉；《辽史》《金史》相关本纪；黑水城西夏文文书与考古报告标明来源路径；该锚点连接编年节点，关于数量、实施和地方后果仍应遵守材料的文体与立场限制。

\paragraph{1189年}\eventanchor{XIA-1189-REGNAL-YEAR}{XIA-1189-REGNAL-YEAR}　西夏以李仁孝在位第50年作为诏令、赋役与外交文书的纪年框架。账本以《宋史》卷485〈夏国传〉；《辽史》《金史》相关本纪；黑水城西夏文文书与考古报告标明来源路径；该锚点连接编年节点，关于数量、实施和地方后果仍应遵守材料的文体与立场限制。

\paragraph{1190年}\eventanchor{XIA-1190-REGNAL-YEAR}{XIA-1190-REGNAL-YEAR}　西夏以李仁孝在位第51年作为诏令、赋役与外交文书的纪年框架。账本以《宋史》卷485〈夏国传〉；《辽史》《金史》相关本纪；黑水城西夏文文书与考古报告标明来源路径；该锚点连接编年节点，关于数量、实施和地方后果仍应遵守材料的文体与立场限制。

\paragraph{1191年}\eventanchor{XIA-1191-REGNAL-YEAR}{XIA-1191-REGNAL-YEAR}　西夏以李仁孝在位第52年作为诏令、赋役与外交文书的纪年框架。账本以《宋史》卷485〈夏国传〉；《辽史》《金史》相关本纪；黑水城西夏文文书与考古报告标明来源路径；该锚点连接编年节点，关于数量、实施和地方后果仍应遵守材料的文体与立场限制。

\paragraph{1192年}\eventanchor{XIA-1192-REGNAL-YEAR}{XIA-1192-REGNAL-YEAR}　西夏以李仁孝在位第53年作为诏令、赋役与外交文书的纪年框架。账本以《宋史》卷485〈夏国传〉；《辽史》《金史》相关本纪；黑水城西夏文文书与考古报告标明来源路径；该锚点连接编年节点，关于数量、实施和地方后果仍应遵守材料的文体与立场限制。

\paragraph{1193年}\eventanchor{XIA-1193-EVENT-071}{XIA-1193-EVENT-071}　李纯佑即位。账本以《宋史》卷485〈夏国传〉；《辽史》《金史》相关本纪；黑水城西夏文文书与考古报告标明来源路径；该锚点连接编年节点，关于数量、实施和地方后果仍应遵守材料的文体与立场限制。

\paragraph{1193年}\eventanchor{XIA-1193-REGNAL-YEAR}{XIA-1193-REGNAL-YEAR}　西夏以李仁孝在位第54年作为诏令、赋役与外交文书的纪年框架。账本以《宋史》卷485〈夏国传〉；《辽史》《金史》相关本纪；黑水城西夏文文书与考古报告标明来源路径；该锚点连接编年节点，关于数量、实施和地方后果仍应遵守材料的文体与立场限制。

\paragraph{1194年}\eventanchor{XIA-1194-REGNAL-YEAR}{XIA-1194-REGNAL-YEAR}　西夏以李纯佑在位第1年作为诏令、赋役与外交文书的纪年框架。账本以《宋史》卷485〈夏国传〉；《辽史》《金史》相关本纪；黑水城西夏文文书与考古报告标明来源路径；该锚点连接编年节点，关于数量、实施和地方后果仍应遵守材料的文体与立场限制。

\paragraph{1195年}\eventanchor{XIA-1195-REGNAL-YEAR}{XIA-1195-REGNAL-YEAR}　西夏以李纯佑在位第2年作为诏令、赋役与外交文书的纪年框架。账本以《宋史》卷485〈夏国传〉；《辽史》《金史》相关本纪；黑水城西夏文文书与考古报告标明来源路径；该锚点连接编年节点，关于数量、实施和地方后果仍应遵守材料的文体与立场限制。

\paragraph{1196年}\eventanchor{XIA-1196-REGNAL-YEAR}{XIA-1196-REGNAL-YEAR}　西夏以李纯佑在位第3年作为诏令、赋役与外交文书的纪年框架。账本以《宋史》卷485〈夏国传〉；《辽史》《金史》相关本纪；黑水城西夏文文书与考古报告标明来源路径；该锚点连接编年节点，关于数量、实施和地方后果仍应遵守材料的文体与立场限制。

\paragraph{1197年}\eventanchor{XIA-1197-REGNAL-YEAR}{XIA-1197-REGNAL-YEAR}　西夏以李纯佑在位第4年作为诏令、赋役与外交文书的纪年框架。账本以《宋史》卷485〈夏国传〉；《辽史》《金史》相关本纪；黑水城西夏文文书与考古报告标明来源路径；该锚点连接编年节点，关于数量、实施和地方后果仍应遵守材料的文体与立场限制。

\paragraph{1198年}\eventanchor{XIA-1198-REGNAL-YEAR}{XIA-1198-REGNAL-YEAR}　西夏以李纯佑在位第5年作为诏令、赋役与外交文书的纪年框架。账本以《宋史》卷485〈夏国传〉；《辽史》《金史》相关本纪；黑水城西夏文文书与考古报告标明来源路径；该锚点连接编年节点，关于数量、实施和地方后果仍应遵守材料的文体与立场限制。

\paragraph{1199年}\eventanchor{XIA-1199-REGNAL-YEAR}{XIA-1199-REGNAL-YEAR}　西夏以李纯佑在位第6年作为诏令、赋役与外交文书的纪年框架。账本以《宋史》卷485〈夏国传〉；《辽史》《金史》相关本纪；黑水城西夏文文书与考古报告标明来源路径；该锚点连接编年节点，关于数量、实施和地方后果仍应遵守材料的文体与立场限制。

\paragraph{1200年}\eventanchor{XIA-1200-REGNAL-YEAR}{XIA-1200-REGNAL-YEAR}　西夏以李纯佑在位第7年作为诏令、赋役与外交文书的纪年框架。账本以《宋史》卷485〈夏国传〉；《辽史》《金史》相关本纪；黑水城西夏文文书与考古报告标明来源路径；该锚点连接编年节点，关于数量、实施和地方后果仍应遵守材料的文体与立场限制。

\paragraph{1201年}\eventanchor{XIA-1201-REGNAL-YEAR}{XIA-1201-REGNAL-YEAR}　西夏以李纯佑在位第8年作为诏令、赋役与外交文书的纪年框架。账本以《宋史》卷485〈夏国传〉；《辽史》《金史》相关本纪；黑水城西夏文文书与考古报告标明来源路径；该锚点连接编年节点，关于数量、实施和地方后果仍应遵守材料的文体与立场限制。

\paragraph{1202年}\eventanchor{XIA-1202-REGNAL-YEAR}{XIA-1202-REGNAL-YEAR}　西夏以李纯佑在位第9年作为诏令、赋役与外交文书的纪年框架。账本以《宋史》卷485〈夏国传〉；《辽史》《金史》相关本纪；黑水城西夏文文书与考古报告标明来源路径；该锚点连接编年节点，关于数量、实施和地方后果仍应遵守材料的文体与立场限制。

\paragraph{1203年}\eventanchor{XIA-1203-REGNAL-YEAR}{XIA-1203-REGNAL-YEAR}　西夏以李纯佑在位第10年作为诏令、赋役与外交文书的纪年框架。账本以《宋史》卷485〈夏国传〉；《辽史》《金史》相关本纪；黑水城西夏文文书与考古报告标明来源路径；该锚点连接编年节点，关于数量、实施和地方后果仍应遵守材料的文体与立场限制。

\paragraph{1204年}\eventanchor{XIA-1204-REGNAL-YEAR}{XIA-1204-REGNAL-YEAR}　西夏以李纯佑在位第11年作为诏令、赋役与外交文书的纪年框架。账本以《宋史》卷485〈夏国传〉；《辽史》《金史》相关本纪；黑水城西夏文文书与考古报告标明来源路径；该锚点连接编年节点，关于数量、实施和地方后果仍应遵守材料的文体与立场限制。

\paragraph{1205年}\eventanchor{XIA-1205-REGNAL-YEAR}{XIA-1205-REGNAL-YEAR}　西夏以李纯佑在位第12年作为诏令、赋役与外交文书的纪年框架。账本以《宋史》卷485〈夏国传〉；《辽史》《金史》相关本纪；黑水城西夏文文书与考古报告标明来源路径；该锚点连接编年节点，关于数量、实施和地方后果仍应遵守材料的文体与立场限制。

\paragraph{1206年}\eventanchor{XIA-1206-EVENT-072}{XIA-1206-EVENT-072}　成吉思汗统一蒙古诸部，西夏北方环境突变。账本以《宋史》卷485〈夏国传〉；《辽史》《金史》相关本纪；黑水城西夏文文书与考古报告标明来源路径；该锚点连接编年节点，关于数量、实施和地方后果仍应遵守材料的文体与立场限制。

\paragraph{1206年}\eventanchor{XIA-1206-REGNAL-YEAR}{XIA-1206-REGNAL-YEAR}　西夏以李纯佑在位第13年作为诏令、赋役与外交文书的纪年框架。账本以《宋史》卷485〈夏国传〉；《辽史》《金史》相关本纪；黑水城西夏文文书与考古报告标明来源路径；该锚点连接编年节点，关于数量、实施和地方后果仍应遵守材料的文体与立场限制。

\paragraph{1207年}\eventanchor{XIA-1207-REGNAL-YEAR}{XIA-1207-REGNAL-YEAR}　西夏以李安全在位第1年作为诏令、赋役与外交文书的纪年框架。账本以《宋史》卷485〈夏国传〉；《辽史》《金史》相关本纪；黑水城西夏文文书与考古报告标明来源路径；该锚点连接编年节点，关于数量、实施和地方后果仍应遵守材料的文体与立场限制。

\paragraph{1208年}\eventanchor{XIA-1208-REGNAL-YEAR}{XIA-1208-REGNAL-YEAR}　西夏以李安全在位第2年作为诏令、赋役与外交文书的纪年框架。账本以《宋史》卷485〈夏国传〉；《辽史》《金史》相关本纪；黑水城西夏文文书与考古报告标明来源路径；该锚点连接编年节点，关于数量、实施和地方后果仍应遵守材料的文体与立场限制。

\paragraph{1209年}\eventanchor{XIA-1209-EVENT-073}{XIA-1209-EVENT-073}　蒙古军攻西夏，西夏接受和议与联姻安排。账本以《宋史》卷485〈夏国传〉；《辽史》《金史》相关本纪；黑水城西夏文文书与考古报告标明来源路径；该锚点连接编年节点，关于数量、实施和地方后果仍应遵守材料的文体与立场限制。

\paragraph{1209年}\eventanchor{XIA-1209-REGNAL-YEAR}{XIA-1209-REGNAL-YEAR}　西夏以李安全在位第3年作为诏令、赋役与外交文书的纪年框架。账本以《宋史》卷485〈夏国传〉；《辽史》《金史》相关本纪；黑水城西夏文文书与考古报告标明来源路径；该锚点连接编年节点，关于数量、实施和地方后果仍应遵守材料的文体与立场限制。

\paragraph{1210年}\eventanchor{XIA-1210-REGNAL-YEAR}{XIA-1210-REGNAL-YEAR}　西夏以李安全在位第4年作为诏令、赋役与外交文书的纪年框架。账本以《宋史》卷485〈夏国传〉；《辽史》《金史》相关本纪；黑水城西夏文文书与考古报告标明来源路径；该锚点连接编年节点，关于数量、实施和地方后果仍应遵守材料的文体与立场限制。

\paragraph{1211年}\eventanchor{XIA-1211-REGNAL-YEAR}{XIA-1211-REGNAL-YEAR}　西夏以李安全在位第5年作为诏令、赋役与外交文书的纪年框架。账本以《宋史》卷485〈夏国传〉；《辽史》《金史》相关本纪；黑水城西夏文文书与考古报告标明来源路径；该锚点连接编年节点，关于数量、实施和地方后果仍应遵守材料的文体与立场限制。

\paragraph{1212年}\eventanchor{XIA-1212-REGNAL-YEAR}{XIA-1212-REGNAL-YEAR}　西夏以李遵顼在位第1年作为诏令、赋役与外交文书的纪年框架。账本以《宋史》卷485〈夏国传〉；《辽史》《金史》相关本纪；黑水城西夏文文书与考古报告标明来源路径；该锚点连接编年节点，关于数量、实施和地方后果仍应遵守材料的文体与立场限制。

\paragraph{1213年}\eventanchor{XIA-1213-REGNAL-YEAR}{XIA-1213-REGNAL-YEAR}　西夏以李遵顼在位第2年作为诏令、赋役与外交文书的纪年框架。账本以《宋史》卷485〈夏国传〉；《辽史》《金史》相关本纪；黑水城西夏文文书与考古报告标明来源路径；该锚点连接编年节点，关于数量、实施和地方后果仍应遵守材料的文体与立场限制。

\paragraph{1214年}\eventanchor{XIA-1214-REGNAL-YEAR}{XIA-1214-REGNAL-YEAR}　西夏以李遵顼在位第3年作为诏令、赋役与外交文书的纪年框架。账本以《宋史》卷485〈夏国传〉；《辽史》《金史》相关本纪；黑水城西夏文文书与考古报告标明来源路径；该锚点连接编年节点，关于数量、实施和地方后果仍应遵守材料的文体与立场限制。

\paragraph{1215年}\eventanchor{XIA-1215-REGNAL-YEAR}{XIA-1215-REGNAL-YEAR}　西夏以李遵顼在位第4年作为诏令、赋役与外交文书的纪年框架。账本以《宋史》卷485〈夏国传〉；《辽史》《金史》相关本纪；黑水城西夏文文书与考古报告标明来源路径；该锚点连接编年节点，关于数量、实施和地方后果仍应遵守材料的文体与立场限制。

\paragraph{1216年}\eventanchor{XIA-1216-REGNAL-YEAR}{XIA-1216-REGNAL-YEAR}　西夏以李遵顼在位第5年作为诏令、赋役与外交文书的纪年框架。账本以《宋史》卷485〈夏国传〉；《辽史》《金史》相关本纪；黑水城西夏文文书与考古报告标明来源路径；该锚点连接编年节点，关于数量、实施和地方后果仍应遵守材料的文体与立场限制。

\paragraph{1217年}\eventanchor{XIA-1217-EVENT-074}{XIA-1217-EVENT-074}　西夏对蒙古军事要求的配合出现紧张。账本以《宋史》卷485〈夏国传〉；《辽史》《金史》相关本纪；黑水城西夏文文书与考古报告标明来源路径；该锚点连接编年节点，关于数量、实施和地方后果仍应遵守材料的文体与立场限制。

\paragraph{1217年}\eventanchor{XIA-1217-REGNAL-YEAR}{XIA-1217-REGNAL-YEAR}　西夏以李遵顼在位第6年作为诏令、赋役与外交文书的纪年框架。账本以《宋史》卷485〈夏国传〉；《辽史》《金史》相关本纪；黑水城西夏文文书与考古报告标明来源路径；该锚点连接编年节点，关于数量、实施和地方后果仍应遵守材料的文体与立场限制。

\paragraph{1218年}\eventanchor{XIA-1218-REGNAL-YEAR}{XIA-1218-REGNAL-YEAR}　西夏以李遵顼在位第7年作为诏令、赋役与外交文书的纪年框架。账本以《宋史》卷485〈夏国传〉；《辽史》《金史》相关本纪；黑水城西夏文文书与考古报告标明来源路径；该锚点连接编年节点，关于数量、实施和地方后果仍应遵守材料的文体与立场限制。

\paragraph{1219年}\eventanchor{XIA-1219-REGNAL-YEAR}{XIA-1219-REGNAL-YEAR}　西夏以李遵顼在位第8年作为诏令、赋役与外交文书的纪年框架。账本以《宋史》卷485〈夏国传〉；《辽史》《金史》相关本纪；黑水城西夏文文书与考古报告标明来源路径；该锚点连接编年节点，关于数量、实施和地方后果仍应遵守材料的文体与立场限制。

\paragraph{1220年}\eventanchor{XIA-1220-REGNAL-YEAR}{XIA-1220-REGNAL-YEAR}　西夏以李遵顼在位第9年作为诏令、赋役与外交文书的纪年框架。账本以《宋史》卷485〈夏国传〉；《辽史》《金史》相关本纪；黑水城西夏文文书与考古报告标明来源路径；该锚点连接编年节点，关于数量、实施和地方后果仍应遵守材料的文体与立场限制。

\paragraph{1221年}\eventanchor{XIA-1221-REGNAL-YEAR}{XIA-1221-REGNAL-YEAR}　西夏以李遵顼在位第10年作为诏令、赋役与外交文书的纪年框架。账本以《宋史》卷485〈夏国传〉；《辽史》《金史》相关本纪；黑水城西夏文文书与考古报告标明来源路径；该锚点连接编年节点，关于数量、实施和地方后果仍应遵守材料的文体与立场限制。

\paragraph{1222年}\eventanchor{XIA-1222-REGNAL-YEAR}{XIA-1222-REGNAL-YEAR}　西夏以李遵顼在位第11年作为诏令、赋役与外交文书的纪年框架。账本以《宋史》卷485〈夏国传〉；《辽史》《金史》相关本纪；黑水城西夏文文书与考古报告标明来源路径；该锚点连接编年节点，关于数量、实施和地方后果仍应遵守材料的文体与立场限制。

\paragraph{1223年}\eventanchor{XIA-1223-EVENT-075}{XIA-1223-EVENT-075}　李德旺即位，西夏王室更替发生在蒙古压力下。账本以《宋史》卷485〈夏国传〉；《辽史》《金史》相关本纪；黑水城西夏文文书与考古报告标明来源路径；该锚点连接编年节点，关于数量、实施和地方后果仍应遵守材料的文体与立场限制。

\paragraph{1223年}\eventanchor{XIA-1223-REGNAL-YEAR}{XIA-1223-REGNAL-YEAR}　西夏以李遵顼在位第12年作为诏令、赋役与外交文书的纪年框架。账本以《宋史》卷485〈夏国传〉；《辽史》《金史》相关本纪；黑水城西夏文文书与考古报告标明来源路径；该锚点连接编年节点，关于数量、实施和地方后果仍应遵守材料的文体与立场限制。

\paragraph{1224年}\eventanchor{XIA-1224-REGNAL-YEAR}{XIA-1224-REGNAL-YEAR}　西夏以李德旺在位第1年作为诏令、赋役与外交文书的纪年框架。账本以《宋史》卷485〈夏国传〉；《辽史》《金史》相关本纪；黑水城西夏文文书与考古报告标明来源路径；该锚点连接编年节点，关于数量、实施和地方后果仍应遵守材料的文体与立场限制。

\paragraph{1225年}\eventanchor{XIA-1225-REGNAL-YEAR}{XIA-1225-REGNAL-YEAR}　西夏以李德旺在位第2年作为诏令、赋役与外交文书的纪年框架。账本以《宋史》卷485〈夏国传〉；《辽史》《金史》相关本纪；黑水城西夏文文书与考古报告标明来源路径；该锚点连接编年节点，关于数量、实施和地方后果仍应遵守材料的文体与立场限制。

\paragraph{1226年}\eventanchor{XIA-1226-EVENT-076}{XIA-1226-EVENT-076}　蒙古军大举攻入西夏。账本以《宋史》卷485〈夏国传〉；《辽史》《金史》相关本纪；黑水城西夏文文书与考古报告标明来源路径；该锚点连接编年节点，关于数量、实施和地方后果仍应遵守材料的文体与立场限制。

\paragraph{1226年}\eventanchor{XIA-1226-REGNAL-YEAR}{XIA-1226-REGNAL-YEAR}　西夏以李德旺在位第3年作为诏令、赋役与外交文书的纪年框架。账本以《宋史》卷485〈夏国传〉；《辽史》《金史》相关本纪；黑水城西夏文文书与考古报告标明来源路径；该锚点连接编年节点，关于数量、实施和地方后果仍应遵守材料的文体与立场限制。

\paragraph{1227年}\eventanchor{XIA-1227-EVENT-077}{XIA-1227-EVENT-077}　西夏灭亡，蒙古接收河西和宁夏区域。账本以《宋史》卷485〈夏国传〉；《辽史》《金史》相关本纪；黑水城西夏文文书与考古报告标明来源路径；该锚点连接编年节点，关于数量、实施和地方后果仍应遵守材料的文体与立场限制。

\paragraph{1227年}\eventanchor{XIA-1227-REGNAL-YEAR}{XIA-1227-REGNAL-YEAR}　西夏以南平王李睍在位第1年作为诏令、赋役与外交文书的纪年框架。账本以《宋史》卷485〈夏国传〉；《辽史》《金史》相关本纪；黑水城西夏文文书与考古报告标明来源路径；该锚点连接编年节点，关于数量、实施和地方后果仍应遵守材料的文体与立场限制。

